% Chapter 7: Physics-Informed Operators and Sobolev-Regularized Losses
% Part I of the book.
% STATUS: book-standard prose (Part I).

\chapter{Physics-Informed Operators and Sobolev-Regularized Losses}
\label{ch:part1-physics-informed-operators-and-sobolev-reg}
\label{sec:pino}
% -----------------------------------------------------------------------------

Everything proved so far in this part is an existence statement: an architecture of a given
shape approximates a continuous operator on a compact input set to within $\varepsilon$
(\chapref{ch:part1-operator-approximation-theory-and-universa}), spectral truncation leaves
an error of known order (\chapref{ch:part1-the-fourier-neural-operator-spectral-theor}), and
the discrete realization converges to a mesh-free continuum map
(\chapref{ch:part1-discretization-invariance-and-continuum-li}). None of this names the
parameters a training procedure returns. Between good approximants existing and a minimizer
of a finite-sample least-squares functional being found sits the question this chapter takes
up: what objective should be minimized, and what does it control?
\index{objective functional!in operator learning}

The standard recipe trains against paired solver output: draw inputs, run the simulation,
regress. It fails in three regimes that recur in computation. A three-dimensional
time-resolved solve makes labels scarce. Observational data---fields recorded below the
resolution at which the surrogate will be queried---carry fewer degrees of freedom than the
target function. And when the physics is varied, labels generated for one viscosity or one
forcing family are not labels for the next. In each case the empirical loss has a large flat
set: any candidate reproducing the recorded samples scores identically whether or not it
satisfies the equation elsewhere. Supervision constrains a surrogate where data
exist and leaves it free where they do not.

The remedy developed here is to make the equation part of the objective.
\index{Physics-Informed Neural Network (PINN)}\emph{Physics-Informed Neural Networks}
(PINNs)~\cite{raissi2019physics} minimize a residual for one instance of a
\emph{partial differential equation} (PDE);
\index{Physics-Informed Neural Operator (PINO)}\emph{Physics-Informed Neural Operators}
(PINO)~\cite{li2021physics} impose the same constraint on a whole parametric family, so that
a single trained map is penalized for violating the equation on every input it is shown. Two
further ingredients are then required: derivatives of a network output with respect to
physical coordinates, cheap enough to appear inside every training step; and a choice of the
norm that measures the residual, since a residual is a distribution, not a function, when the
trial field is too rough for the differential expression. Both come from machinery already
assembled: the Sobolev scale and duality pairing of
\chapref{ch:part1-function-spaces-duality-and-sobolev-regula}, the weak formulations of
\chapref{ch:part1-partial-differential-equations-as-solution}, and the multiplier calculus of
\chapref{ch:part1-the-fourier-neural-operator-spectral-theor}.

The chapter proceeds in three moves: the composite loss and what its terms control; how the
residual is differentiated on a grid, and where that computation stops being honest; and what
derivative-based metrics buy. The accounting is explicit throughout: physics enters as a
penalty, penalties cost derivative evaluations, and no choice of weights converts a
regularized least-squares problem into a convergence proof.

% -----------------------------------------------------------------------------
\section{The Variational Loss Functional}
\index{variational loss functional}

Let $\mathcal{N}(\,\cdot\,;a)$\index{$\mathcal{N}$@$\mathcal{N}$, differential expression}
denote a nonlinear differential expression depending on an input field $a$, and let
$\mathcal{B}$\index{$\mathcal{B}$@$\mathcal{B}$, boundary operator} be a boundary or initial
operator. The parametric problems collected in
\chapref{ch:part1-partial-differential-equations-as-solution} all fit the form
\begin{equation}
  \label{eq:ch07-bvp}
  \begin{cases}
    \mathcal{N}(u; a) = 0, & x \in \Omega, \\
    \mathcal{B}(u; a) = 0, & x \in \partial\Omega.
  \end{cases}
\end{equation}
For a stationary problem $\mathcal{N}$ is the elliptic expression $-\nabla\cdot(a\nabla u)-f$
of Darcy flow; for an evolutionary one the unknown $u$ is read as a trajectory in the
Bochner space $L^2(0,T;\mathcal{U})$ introduced in
\chapref{ch:part1-function-spaces-duality-and-sobolev-regula} and $\mathcal{N}$ contains
$\partial_t$. The exact solution map is $\mathcal{G}_{\text{true}}:a\mapsto u$, the surrogate
is $\mathcal{G}_\theta$. Inputs are drawn from a probability measure
\index{$\mu$@$\mu\in\mathcal{P}(\mathcal{A})$, training measure}$\mu\in\mathcal{P}(\mathcal{A})$, the same measure whose
support supplies the compact sets of the universality
theorems. All expectations below are population quantities; stochastic training replaces
each by a minibatch average, which is the only place where discretization enters the
objective.

The composite training loss for a parameter vector $\theta$ is
\begin{equation}
  \label{eq:pino_loss}
  \mathcal{L}_{\text{PINO}}(\theta) = \lambda_{\text{data}} \mathcal{L}_{\text{data}}(\theta) + \lambda_{\text{pde}} \mathcal{L}_{\text{pde}}(\theta) + \lambda_{\text{bc}} \mathcal{L}_{\text{bc}}(\theta),
\end{equation}
\index{PDE residual penalty}\index{loss weighting!$\lambda$ coefficients}where
\begin{align}
  \mathcal{L}_{\text{data}}(\theta) &= \mathbb{E}_{a \sim \mu} \left[ \frac{\|\mathcal{G}_\theta(a) - \mathcal{G}_{\text{true}}(a)\|_{L^2(\Omega)}^2}{\|\mathcal{G}_{\text{true}}(a)\|_{L^2(\Omega)}^2} \right], \label{eq:ch07-data-loss}\\
  \mathcal{L}_{\text{pde}}(\theta) &= \mathbb{E}_{a \sim \mu} \left[ \|\mathcal{N}(\mathcal{G}_\theta(a); a)\|_{X(\Omega)}^2 \right], \label{eq:ch07-pde-loss}\\
  \mathcal{L}_{\text{bc}}(\theta) &= \mathbb{E}_{a \sim \mu} \left[ \|\mathcal{B}(\mathcal{G}_\theta(a); a)\|_{Y(\partial\Omega)}^2 \right]. \label{eq:ch07-bc-loss}
\end{align}
Set $\lambda_{\text{pde}}=\lambda_{\text{bc}}=0$ in \eqref{eq:pino_loss} and only
\eqref{eq:ch07-data-loss} survives: that is supervised regression on paired simulations, the
procedure the previous chapters analyzed. The two added terms
\eqref{eq:ch07-pde-loss}--\eqref{eq:ch07-bc-loss} measure, respectively, how much the
predicted field violates the equation inside $\Omega$ and how much it violates the imposed
trace on $\partial\Omega$.

Two features of \eqref{eq:ch07-data-loss}--\eqref{eq:ch07-bc-loss} repay attention. Only the
data term is normalized by the energy of its target, so it is scale-free and dimensionless;
the residual and boundary terms carry the physical units of $\mathcal{N}$ and
$\mathcal{B}$. Consequently $\lambda_{\text{pde}}$ and $\lambda_{\text{bc}}$ are dimensional
quantities whose usable size depends on the nondimensionalization chosen for the problem,
and balancing the three terms is a hyperparameter decision rather than a consequence of the
formulation. Next, the terms are not independent constraints but penalties: what
physics-informed training really asserts is a conditional. If a parameter $\theta^{*}$
annihilates all three terms, then $\mathcal{G}_{\theta^{*}}(a)=\mathcal{G}_{\text{true}}(a)$
for $\mu$-almost every $a$, \emph{provided} the problem \eqref{eq:ch07-bvp} has a unique
solution in the class containing both fields. Well-posedness---the Lax--Milgram argument of
\chapref{ch:part1-partial-differential-equations-as-solution}---is therefore a hypothesis of
the method, not background decoration. What the residual term changes about data
requirements is visible in this statement: instead of requiring a full solution field as a
label, an input $a$ together with the equation supplies a target, and the two can be
combined at different resolutions, as in the coarse-data, fine-residual training reported
in~\cite{li2021physics}.

The choice of the space $X(\Omega)$ in \eqref{eq:ch07-pde-loss} is where the functional
analysis does work. With $X=L^2(\Omega)$ the residual is evaluated in the \emph{strong form},
pointwise on the grid, which demands that the trial field possess as many classical
derivatives as $\mathcal{N}$ has. With $X=H^{-s}(\Omega)$, $s>0$, the residual is measured in
the \emph{weak form}: for a trial field $v=\mathcal{G}_\theta(a)$ whose residual is denoted
$r_v$, the norm is the dual one,
$\|r_v\|_{H^{-s}}=\sup_{\varphi\neq0}|\langle r_v,\varphi\rangle|/\|\varphi\|_{H^{s}}$, and
integration by parts transfers derivatives from $v$ onto the smooth test function
$\varphi$. For a second-order divergence-form operator the
strong residual needs two derivatives of $v$ while the weak residual needs one. By Parseval's
identity the same distinction is a weighting of modal coefficients,
$\|r_v\|_{H^{-s}}^{2}=\sum_{k}(1+|k|^{2})^{-s}|\hat{r}_v(k)|^{2}$---a relaxation of
attention on fine scales, in direct contrast with the sharpening that the Sobolev-norm
penalty of the third section performs.

For coercive problems the weak penalty is more than a convention: it controls the error. Take
the Darcy problem under uniform ellipticity, with $v\in H^1_0(\Omega)$ a trial field of the
correct trace and $r_v$ its weak residual,
$\langle r_v,w\rangle:=\int_\Omega a\nabla v\cdot\nabla w\,dx-\langle f,w\rangle_{H^{-1},H^1_0}$.
Subtracting the weak form of the exact solution from that of $v$, testing with
$w=v-\mathcal{G}_{\text{true}}(a)$, and using coercivity
$B_a(e,e)\ge a_{\min}\|\nabla e\|_{L^2}^2$ gives the a priori estimate
\begin{equation}
  \label{eq:ch07-apriori}
  \big\|\nabla(v-\mathcal{G}_{\text{true}}(a))\big\|_{L^2(\Omega)}
  \;\le\; \frac{1}{a_{\min}}\,\|r_v\|_{H^{-1}(\Omega)}.
\end{equation}
Here $\|\cdot\|_{H^{-1}(\Omega)}$ is the dual of the energy norm
$\|\nabla\cdot\|_{L^2(\Omega)}$ on $H^1_0(\Omega)$; measured against the full $H^1$ norm, a
Poincar{\'e} constant joins the right-hand side, as in
\chapref{ch:part1-partial-differential-equations-as-solution}. This is the precise sense in
which minimizing \eqref{eq:ch07-pde-loss} substitutes for
labeling: the estimate \eqref{eq:ch07-apriori} says that a small weak residual plus a correct
trace forces a small energy-norm error. Read it
also as a warning, because the estimate is conditional on coercivity and degrades as
$a_{\min}\to0$; where no such estimate is available---the inviscid limit, formation of
shocks---a small residual is not evidence of a small error. It is worth recording that the
empirical ranking does not follow the analysis: in the Darcy experiments reported in~\cite{li2021physics},
the strong-form residual trained better than the weak-form one, which
suggests that the test space is an implementation choice entangled with how derivatives are
computed, taken up next.

The boundary term deserves one further comment. On the torus $\mathbb{T}^d$ there is no
$\partial\Omega$, so \eqref{eq:ch07-bc-loss} vanishes identically and the vorticity-form
Navier--Stokes problem of \chapref{ch:part1-partial-differential-equations-as-solution} is
trained on data and residual alone. Where a boundary does exist, \eqref{eq:ch07-bc-loss} is
still a soft penalty. The alternative is a hard constraint---compose the surrogate with a
lifting of the boundary data so that the trace is satisfied by construction---which removes
one weight from the hyperparameter list but requires the trace operator to be meaningful for
the output regularity at hand, that is $v\in H^{s}(\Omega)$ with $s>1/2$.

% -----------------------------------------------------------------------------
\section{Spectral Differentiation in Operator Learning}
\index{spectral differentiation}\index{pseudo-spectral method}

In a PINN~\cite{raissi2019physics} the network is itself the solution ansatz,
$u_\theta:x\mapsto u(x)$, evaluated at scattered collocation points, and the spatial
derivatives $\partial_x u$ and $\partial_{xx}u$ in the residual are obtained by reverse-mode
automatic differentiation. Nothing is approximated there in the numerical-analysis sense:
the derivative is exact for the model, since the model is a differentiable function of $x$.
The price is structural. A second derivative requires a second backward pass through the
computational graph, time and memory grow with depth times the number of collocation points,
and each point is evaluated independently, so representing a field means representing its
samples.

A neural operator is organized the other way round. Its output is a field on a discretization
$\mathcal{M}_N$, and its internal spectral layer---a learned Fourier multiplier, as
\chapref{ch:part1-the-fourier-neural-operator-spectral-theor} constructs it---is a global
operation on that field, defined only for the whole grid at once. Querying such a network at
a single coordinate $x$ means running the transforms; forcing it into a form on which
automatic differentiation yields second $x$-derivatives is possible but forfeits exactly the
structure that made the architecture cheap. When the output is a discrete field on a periodic
mesh $\mathbb{T}^d$, spatial derivatives can instead be evaluated directly in Fourier space:
\begin{equation}
  \label{eq:spectral_diff}
  \frac{\partial^m \mathcal{G}_\theta(a)}{\partial x_j^m} = \mathcal{F}^{-1} \left( (2\pi i k_j)^m \cdot \mathcal{F}\big( \mathcal{G}_\theta(a) \big)(k) \right).
\end{equation}
The identity is one line. For $N$ equispaced nodes of $[0,1)^d$ with discrete Fourier
coefficients $\hat v(k)$ of a grid function $v$, differentiation of the trigonometric
interpolant acts diagonally,
\begin{equation}
  \label{eq:ch07-diff-prop}
  \big(\mathcal{F}(\partial_{x_j}^{m} v)\big)(k) = (2\pi i k_j)^{m}\,\hat v(k),
  \qquad k\in\mathbb{Z}^{d},
\end{equation}
and \eqref{eq:spectral_diff} is \eqref{eq:ch07-diff-prop} followed by the inverse transform.
\index{Fourier multiplier!differentiation as}

Two consequences make this the right tool for a physics-informed operator. First,
\eqref{eq:ch07-diff-prop} is a Fourier multiplier with symbol $(2\pi i k_j)^m$: the very
primitive the architecture already applies, with the learned tensor $R_\phi(k)$ replaced by a
fixed one. Differentiating a field costs two traversals of the
\emph{fast Fourier transform} (FFT) pipeline, $\mathcal{O}(N\log N)$ operations for the
entire field at once, rather than a per-point differentiation of a network graph. Second,
\eqref{eq:spectral_diff} is exact on band-limited inputs and, for smooth ones, converges
faster than any power of the mesh width: if $v$ extends analytically to a complex
neighborhood of the torus, its coefficients decay geometrically, so truncating at $k_{\max}$
leaves an error in the $m$-th derivative of the same exponential order, whereas a finite
difference of order $p$ carries a truncation error $\mathcal{O}(h^{p})$ that shrinks only
algebraically~\cite{leveque2002finite}. \index{exponential convergence!of spectral differentiation}

Both statements need qualifying, and the qualification matters more than the advantage. What
\eqref{eq:spectral_diff} differentiates is the interpolant of the stored samples, not the
continuum field they came from. Where the two differ, the discrepancy is amplified rather
than damped, because the symbol in \eqref{eq:ch07-diff-prop} has magnitude
$(2\pi|k|)^{m}$: an error component at the highest resolvable wavenumber $k\approx N/2$ is
multiplied by $(\pi N)^{m}$ in the $m$-th derivative, so grid-scale noise or truncation
content dominates a residual built this way. For even $N$ the Nyquist coefficient is not
resolved by the samples at all---two interpolants agree at the nodes and disagree in their
derivative---and it is conventionally discarded before an odd-order
differentiation~\index{aliasing!in spectral derivatives}\index{Nyquist mode}. Nor is a
nonlinear residual free of the same difficulty.
\begin{example}[pseudo-spectral Burgers residual]
For $\mathcal{N}(u) = \partial_t u + u\,\partial_x u - \nu\,\partial_{xx}u$ on $\mathbb{T}$,
each linear term comes from \eqref{eq:spectral_diff}, while the product $u\,\partial_x u$ is
formed pointwise on the grid and, if it is ever transformed again, contains wavenumbers
beyond the cutoff that fold back onto resolved modes. Classical dealiasing (zero-padding or
the two-thirds rule) suppresses that \emph{aliasing}; a PINO implementation inherits the
requirement because the residual is recomputed every step. The same restriction is recorded
in~\cite{li2021physics}: spectral differentiation presumes smoothness and a uniform grid, and
derivative error is amplified downstream in the predicted solution. For fields with jumps,
such as the vanishing-viscosity Burgers limit discussed in
\chapref{ch:part1-partial-differential-equations-as-solution}, no nodal differentiation rule
is accurate---the Gibbs phenomenon identified as an open problem in
\chapref{ch:part1-open-theoretical-problems-and-mathematical} is precisely this failure.
\end{example}

Finally, note where the differentiation is \emph{not} taken. In
\eqref{eq:spectral_diff} the derivatives are with respect to the physical coordinates $x$, at
fixed $\theta$; the training gradient still requires derivatives with respect to $\theta$,
computed by backpropagation through the field-producing network. Because the operator in
\eqref{eq:ch07-diff-prop} is linear, its adjoint is the multiplier with the conjugated
symbol, so the backward pass through a spectral derivative is another pair of FFTs rather
than a second-order differentiation of the network graph. The practical cost is therefore
memory, not algebra: intermediates of the forward pass must be retained so that the residual
branch can be differentiated, and the residual must be assembled at every grid point of every
sample in every step.

% -----------------------------------------------------------------------------
\section{Sobolev \texorpdfstring{$H^1$ and $H^2$}{H1 and H2} Norm Regularization}
\index{Sobolev space!$H^1$ regularization}\index{Sobolev space!$H^2$ regularization}

Training neural operators solely with the relative $L^2$ loss \eqref{eq:ch07-data-loss} can
lead to \emph{spectral bias}\index{spectral bias}, where the model captures low-frequency
energy but exhibits phase errors or high-frequency oscillations in gradients. The mechanism
is visible in Fourier coefficients. Write the prediction error as
$e_\theta(a)=\mathcal{G}_\theta(a)-u$, where $u=\mathcal{G}_{\text{true}}(a)$ is the label
field. Parseval's identity gives
$\|e_\theta\|_{L^2}^{2}=\sum_k|\hat e_\theta(k)|^{2}$ while
$\|\partial_{x_j}e_\theta\|_{L^2}^{2}=\sum_k(2\pi k_j)^2|\hat e_\theta(k)|^{2}$, so a defect
at wavenumber $k$ counts $(2\pi|k|)^2$ times more heavily in the derivative metric than in
the $L^2$ metric. Since the target's own spectrum decays with $|k|$
(cf.\ the coefficient bound in \chapref{ch:part1-the-fourier-neural-operator-spectral-theor}),
residual error is concentrated exactly where relative $L^2$ error is nearly blind, and a
first-order optimizer spends its early iterations on the high-energy modes.

Regularizing the objective with higher Sobolev norms enforces derivative fidelity:
\begin{equation}
  \label{eq:ch07-sobolev-loss}
  \mathcal{L}_{H^s}(\theta) = \mathbb{E}_{a \sim \mu} \left[ \frac{\|\mathcal{G}_\theta(a) - u\|_{L^2}^2}{\|u\|_{L^2}^2} + \beta \sum_{j=1}^d \frac{\|\partial_{x_j} \mathcal{G}_\theta(a) - \partial_{x_j} u\|_{L^2}^2}{\|\partial_{x_j} u\|_{L^2}^2} \right].
\end{equation}
The sum runs over the $d$ coordinate
directions, and $\beta$\index{$\beta$@$\beta$, derivative-tradeoff weight} balances the two
quotients. In spite of the subscript, \eqref{eq:ch07-sobolev-loss} controls values and first
derivatives only, that is the inhomogeneous $H^1$ norm; the $H^2$ case named in the section
title appends the second multi-index terms, which in Fourier variables means weighting modal
errors by $(1+|k|^{2})^{2}$ rather than $(1+|k|^{2})^{1}$. Read this way, \eqref{eq:ch07-sobolev-loss}
and the negative-norm residual penalty of \eqref{eq:ch07-pde-loss} are the same device on
opposite sides of the Sobolev scale\index{Sobolev space!as objective weighting}: a positive
index sharpens attention on fine structure, a negative index relaxes it.

\index{Sobolev embedding theorem!and $L^\infty$ control}What the sharpening buys is a change
of metric, and one metric inequality. The Sobolev embedding theorem
(\chapref{ch:part1-function-spaces-duality-and-sobolev-regula}; see
also~\cite{evans2010partial,brezis2011functional}) states that $H^{s}(\Omega)\hookrightarrow
C^{k}(\overline{\Omega})$ continuously when $s>d/2+k$, so with $k=0$ an error $e$ obeys
$\|e\|_{L^\infty}\le C_s\|e\|_{H^{s}}$. Controlling $\|\cdot\|_{H^s}$ with $s>d/2$ therefore
excludes localized pointwise spikes: no defect can be large somewhere without being large in
the norm the objective measures. The statement is an implication between metrics, not an
algorithmic guarantee---nothing in \eqref{eq:ch07-sobolev-loss} forces the optimizer to make
that norm small, and in dimension $d=3$ the hypothesis $s>d/2$ requires two derivatives,
which is why the $H^2$ variant rather than the $H^1$ one is the natural target for
three-dimensional fields. It is also why such losses are attractive independently of any
physics term: the quantities practitioners read off a simulation are derivatives of the
state---the Darcy flux $-a\nabla u$, the vorticity $\nabla\times u$, the wall shear
stress---so an $L^2$-accurate field may still be a poor model of the diagnostic. The
companion claim, that derivative regularization preserves conservation properties, should be
read in the same register: it suppresses non-physical oscillation in gradients, which is
consistent with better behaved balance laws, but it is not a proof that a trained network
satisfies any conservation identity.

\index{discretization invariance!of the loss functional}A second consistency check comes from
\chapref{ch:part1-discretization-invariance-and-continuum-li}. Each quotient in
\eqref{eq:ch07-sobolev-loss} approximates a ratio of integrals, so on $N$ nodes the sums must
carry the cell volume $h^{d}$ to be quadratures at all; the relative normalization then cancels
the remaining mesh factor, which is why the objective is a mesh-consistent quantity whereas an
unnormalized discrete seminorm grows with refinement. The normalization has its own failure
mode: when $\|\partial_{x_j}u\|_{L^2}$ is small---a label that is nearly constant along a
coordinate, as for shear-driven or symmetry-reduced data---the corresponding term is
ill-conditioned and its gradient is dominated by the surrogate's own derivative fluctuation.
Normalizing over a batch instead of per sample repairs the scaling but makes an individual
sample's contribution depend on which other samples were drawn.

\index{no free lunch!in physics-informed regularization}It remains to state what the two
regularization devices cost, since neither is free. A Sobolev-metric loss requires
differentiating the label as well as the prediction; for fixed grid-function labels those
derivatives can be computed once offline, but the prediction derivatives recur at every
step. A physics-informed loss requires assembling \eqref{eq:ch07-bvp} on the surrogate's
output at every step, adding spectral derivatives, nonlinear pointwise products, and a
backward pass through them, and adding two further hyperparameters to
\eqref{eq:pino_loss}. Neither penalty alters the approximation classes of
\chapref{ch:part1-operator-approximation-theory-and-universa}: they change which element of
the class is found, not which elements exist. Nor do they immunize training against the
equation being wrong---when $\mathcal{N}$ encodes a closure or a solver's discretization
rather than the physics of interest, the residual term biases the fit away from the data that
was generated by that same model. The reported benefits are correspondingly empirical and
specific: PINO trained with both terms is reported to need fewer labels than a purely
data-driven \emph{Fourier Neural Operator} (FNO), to generalize better at resolutions
above the training one, and on transient
and Kolmogorov flows to reduce average relative error by roughly $7\%$, while remaining
trainable in a data-free regime where a PINN is reported to fail on the optimization
landscape of the PDE constraint~\cite{li2021physics,raissi2019physics}. Those are benchmark
outcomes, not theorems, and they are the honest boundary of what this chapter can claim.

% -----------------------------------------------------------------------------

What this chapter leaves unresolved is comparative. It treated the objective as if it were
architecture-independent, and \eqref{eq:pino_loss} is not: the residual term is cheap when
\eqref{eq:spectral_diff} applies, which requires a uniform periodic grid and a field smooth
enough that the multiplier's $(2\pi|k|)^m$ amplification is harmless; it must be rebuilt from
graph or mesh-based difference operators on the unstructured discretizations that the
geometry-informed architectures use; and for an architecture whose input is a fixed sensor
vector, the data term and the residual term are not measured on comparable discretizations at
all. Likewise the derivative penalties of \eqref{eq:ch07-sobolev-loss} presuppose that
$\partial_{x_j}\mathcal{G}_\theta(a)$ has a well-defined discrete meaning. \chapref{ch:part1-comparative-theoretical-taxonomy-of-archit}
assembles these axes into a single comparison---input format, coordinate query, kernel
parameterization, cost per layer, discretization invariance, and support for physics-residual
evaluation---across DeepONet, the generic neural operator, the FNO, and the geometry-informed
neural operator. Reading it after this chapter is worthwhile precisely because the two
decisions it tabulates, how to represent a map and what to minimize, are one decision made
twice.
